%%%
%%% Implementation section, 
%%% how to use available Agents to solve the problem, case studies, etc.
%%%
\section{Implementation and Case Studies} %main part/hauptteil
\label{sec:implementation}
% Verschiedene Ansätze vorbringen
% App 1: BoxCipher / CryptoBox
% App 2: Route Cipher
% App 3: Bitcoin attacks
% TIP: Report the earlier runs (1 complete, several partial and undocumented) as a pilot study. They informed the AGENTS.md and the experimental protocol. Runeson & Höst (2009) explicitly allow pilot cases, so they need not be discarded.
\subsection{Experimental Setup} % (Kürzen wegen section "Development of ...")
% TIP: Write the experimental protocol (initial prompt, follow-up rules, done/abort criteria, starting state, data collection) down BEFORE the first run. It then becomes this subsection almost directly.
% Describe CTO project context (repo structure, how new apps are added, etc.)
% What CTO apps are to be implemented (and why)? (How many apps should we do? 2?) State differences between apps
% TIP: Choose app 2 clearly different from app 1, ideally more complex. App 2 is the replication case: does the setup effect from app 1 also occur for a different task?
% TIP: Make sure the task's app does not already exist in the repo and is not publicly available in a form the model may have seen in training. Otherwise mention it as a limitation.
% How are the apps described to the AI?
% TIP: Fixed initial prompt per app, identical for all agents and both configurations. Put both prompts verbatim into the appendix.
% Describe the specs given to the agents in each run (one shot, only /init, tools, skills, etc. (not all combinations, but maybe 2-3 different approaches))
% TIP: Final design (8 runs), 2 configurations = baseline (branch ma-base = main minus AGENTS.md and CLAUDE.md, i.e. "no context file" condition of Gloaguen et al.; main has no skills or MCP configs) vs. setup (branch ma-setup, see 04development.tex). Both based on main 6ed00c72; never work on these branches themselves, they are the fixed starting points.
%   - App 1: 3 agents (Claude Code, Codex, OpenCode with local model) x 2 configurations = 6 runs
%   - App 2: Claude Code only x 2 configurations = 2 runs (replication)
%   Extension if time permits: Codex and OpenCode on app 2 as well (= 12 runs); otherwise future work.
% TIP: Run order, so that an early stop still leaves a usable subset:
%   1. App 1, Claude Code, setup     2. App 1, Claude Code, baseline
%   3. App 1, Codex, setup           4. App 1, OpenCode, setup
%   5. App 2, Claude Code, setup     6. App 2, Claude Code, baseline
%   7. App 1, Codex, baseline        8. App 1, OpenCode, baseline
% TIP: Check that every agent actually loads the setup: Claude Code may read CLAUDE.md instead of AGENTS.md (if so, add a CLAUDE.md containing only "@AGENTS.md"). Skill folders differ per agent (.claude/skills, .agents/skills, .opencode/skills). Verify in each transcript that the skills were used.
% TIP: Local model for OpenCode: test tool calling in a short trial run before the real run. Document model, quantization, runtime (e.g., llama.cpp version), context size, and hardware (RTX 3060 12 GB, Ryzen 5 2600, 46 GB RAM).
% NOTE: Local model setup for OpenCode (state 2026-11-10, values read from the machine; update if anything changes before the runs)
%   Hardware:
%   - CPU: AMD Ryzen 5 2600 (6 cores / 12 threads); RAM: 48 GB DDR4 (46 GiB usable), 4 modules = 2x 16 GB (Corsair CMK32GX4M2E3200C16) + 2x 8 GB (Corsair CMK16GX4M2B3000C15), dual channel (16 + 8 GB per channel), running at 3000 MT/s -> theoretical bandwidth ~48 GB/s (relevant because the offloaded experts are read from RAM)
%   - GPU: NVIDIA GeForce RTX 3060, 12 GB VRAM, driver 620.70.06; the GPU also drives the desktop, so somewhat less than 12 GB is free
%   - OS: CachyOS Linux (rolling), kernel 7.3.3-1-cachyos-bore
%   Model:
%   - Qwen3.6-35B-A3B (Alibaba Qwen team), mixture-of-experts: 35B total, ~3B active parameters per token (TODO: verify license and exact numbers on the model card)
%   - GGUF, quantization Q4_K_M, from lmstudio-community (21.2 GB file), downloaded via LM Studio:
%       ~/.lmstudio/models/lmstudio-community/qwen/qwen3.6-35b-a3b/Qwen3.6-35B-A3B-Q4_K_M.gguf
%   - Why this model: as MoE, its expert weights can live in system RAM while attention, shared weights and the KV cache stay on the GPU -> usable speed on 12 GB VRAM.
%   - Considered and rejected: Devstral Small 2 24B (Mistral, EU-based; dense, Q4_K_M = 14.3 GB, does not fit 12 GB -> future work), Qwen3.6-27B / Qwen3.8-27B (dense, too large for a useful context), Bonsai 27B (1-bit version of Qwen3.6-27B; vendor model card reports a clear drop on agentic/tool-calling benchmarks, needs a llama.cpp fork), Gemma 4 E4B (reported too weak for multi-step tool calling). Backup if Qwen fails the tool-calling trial: Gemma 4 26B A4B (also MoE).
%   Runtime:
%   - llama.cpp llama-server, version 0.9.0-dev, build 12211, commit 7d7f02875 (2026-11-08); built from source with CUDA 13.4 (nvcc V13.4.92) and GCC 16.2.1
%   - Chosen over LM Studio / Ollama: per-layer expert offloading (--n-cpu-moe), explicit chat template (--jinja), pinned commit and a full command line for reproducibility
%   - TODO: freeze llama.cpp and CUDA (no rebuild, cuda in pacman IgnorePkg) until all local runs are done
%   Server command (latest):
%       llama-server -m ~/.lmstudio/models/lmstudio-community/qwen/qwen3.6-35b-a3b/Qwen3.6-35B-A3B-Q4_K_M.gguf -ngl 99 --n-cpu-moe 26 -c 32768 --jinja --port 8080
%   - -ngl 99 = all layers on the GPU; --n-cpu-moe 26 = expert weights of the first 26 layers stay in system RAM; -c 32768 = context window in tokens; --jinja = chat template from the GGUF (needed for tool calls); --port 8080 = OpenAI-compatible API at http://127.0.0.1:8080/v1
%   - Observed: ~30 tokens/s generation (TODO: measure properly, incl. prompt processing speed and speed with a nearly full context)
%   - The server created 4 slots (--parallel auto); the log reports 32768 tokens context per slot, i.e. each request can use the full window, while the slots share one cache. Kept the default (OpenCode works mostly sequentially); revisit only if the log shows repeated full prompt reprocessing
%   - Sampling parameters seen in a request (/slots endpoint): temperature 1.0, top_k 20, top_p 0.95, min_p 0.05. TODO: check whether they come from OpenCode or the GGUF defaults; keep identical in both arms
%   - Thinking: the chat template in the GGUF supports enable_thinking (on/off) and preserve_thinking, but no reasoning_effort levels (so no "medium"; only a hard token cap via --reasoning-budget). Thinking is the model default. An observed request with an empty <think></think> block (1,221 tokens) was probably an OpenCode side request (e.g. title generation), not the main session.
%     Tested 2026-11-10: OpenCode's thinking variants (default/low/medium/high) only send "reasoning_effort" (default sends nothing); llama-server + Qwen template ignore it -> identical output for all variants. The per-request "thinking_budget_tokens" (llama.cpp web UI) and --reasoning-budget N are a hard server-side cutoff (model is interrupted, not trained to plan for it), not an effort level. OpenCode also sends max_completion_tokens 8192 (from limit.output), which includes thinking tokens. -> Use OpenCode variant "default" and document it; set any budget server-side only.
%     Plan: --reasoning on --chat-template-kwargs '{"preserve_thinking":true}' (preserve_thinking reportedly recommended for agents on the model card, TODO: verify on the official card). TODO: in the trial, check that thinking blocks appear and whether 32k context fills too fast (then drop preserve_thinking first). Record the setting, keep identical in both arms.
%     Literature for the discussion: thinking budget improves math/coding/STEM benchmarks (Qwen3 Technical Report, Yang et al. 2025, arXiv:2505.09388); but overthinking lowers success in agentic SWE-bench tasks (Cuadron et al. 2025, "The Danger of Overthinking", arXiv:2502.08235) and longer reasoning can lower accuracy (Gema et al. 2025, "Inverse Scaling in Test-Time Compute", TMLR, arXiv:2507.14417). Effort levels of Claude Code / Codex are not comparable to Qwen's on/off switch -> record each agent's setting, justify as vendor default for agentic use.
%   - TODO: check whether 32k context is enough for the setup arm (skills + MCP tool definitions); monitor via curl localhost:8080/slots (field n_prompt_tokens)
%   OpenCode connection:
%   - OpenCode v2.2.7; provider defined globally in ~/.config/opencode/opencode.jsonc, i.e. outside the repo and therefore identical for both arms:
%       provider "llamacpp", package "@opencode/ai/providers/openai-compatible", settings.baseURL "http://127.0.0.1:8080/v1", model "qwen3.6-35b-a3b" with limit context 65536 / output 8192
%   - Start with: opencode -m llamacpp/qwen3.6-35b-a3b (not the lmstudio/ entry, which points to LM Studio on port 1234)
%   - The project-level opencode.json in ma-setup only adds the MCP servers (chakra-ui, playwright); ma-base has no OpenCode config
%   - Verified 2026-11-10: a test prompt via "opencode run" was answered by the llama-server model. TODO: tool-calling trial on a scratch clone; check that the context limit triggers compaction
%   - TODO: put the full opencode.jsonc and the server command into the appendix
% Introduce selection of agents (and why)? (criteria: cost, performance, ease of use, etc.)
% Describe environment setup (i.e., terminal based, desktop app or plugin in VSC) (e.g., how to connect agents to the CTO codebase, how to evaluate their performance, etc.)
% TIP: Identical starting state for every run:
%   - fresh single-branch clone per run, outside ~/Projects/cto/ (git clone --single-branch --branch ma-setup ... ~/ab-runs/setup-1, then git switch -c ab/setup-1 && npm ci; same with ma-base). Never use linked worktrees: in a worktree without .claude/skills, Claude Code loads the main checkout's skills. Optional template folder per arm, copied with cp -a; never start an agent inside the template.
%   - same AGENTS.md (record git hash)
%   - disable user-level config and memory: ~/.claude/CLAUDE.md, Claude Code auto-memory, ~/.codex/AGENTS.md, global MCP servers, ~/.agents/skills/find-skills; record remaining user-level skills, plugins, claude.ai connectors, Claude in Chrome
%   - setup arm only: install Playwright browser once (npx -y @playwright/mcp@0.0.83 install-browser chrome-for-testing) and approve the project MCP servers
%   - same agent, model, prompt and plan-approval policy in both arms
%   - record model, agent version, and reasoning effort per run
% Describe observations that will be made (e.g., (no. of) prompts, no. of iterations, time taken (elapsed time), autonomy, failures, interventions needed, etc.)
% TIP: Rules for follow-up prompts. Classify every intervention:
%   - Report: only describe what is wrong ("build fails with X", "button does nothing")
%   - Hint: point to a place ("look at how app Y does it")
%   - Instruction: say exactly what to do, or write code yourself
%   Stay at "Report" as long as possible. Number of interventions per category and agent = central comparable result.
% TIP: Define "done" and "abort" in advance:
%   - done = checklist (builds, lint passes, required features work, follows app structure)
%   - abort = fixed budget (e.g., 90 min or 25 prompts), after which the run counts as failed
%   Otherwise agents get stopped at different points unintentionally.
% TIP: Data collection, right after each run:
%   - save transcripts: Claude Code ~/.claude/projects/ (JSONL), Codex ~/.codex/sessions/, OpenCode export
%   - commit after every agent turn so the git history shows how the code evolved
%   - token usage / cost can be extracted from the transcripts later
%   - measure: cto-app-check rubric, number of corrective follow-ups, tokens, time, lint and test results. Keep the rubric outside both clones, never let it into a baseline run.
%   - blind rating: git diff ma-setup...ab/setup-1 > ~/ab-ratings/run-07.diff (neutral names), mapping name -> run in a separate file
% TIP: Keep a log table during the run: one row per prompt with time, intervention category, what went wrong, short note. Much faster to analyse than re-reading transcripts.
\subsection{Case Study: Claude Code} %Or: Case Study: App 1; Case Study: App 2; Case Study: App 3
% In-depth description of different approaches (e.g., with/without tools, with/without skills, etc.) 
% 1. prompt/spec given to the agent
% 2. step by step description of what the agent did
% 3. observations (where it got stuck, where intervention was needed)
% 4. results (code snippets, screenshots, etc.)
\subsection{Case Study: Codex}
% same as above (not too deep/detailed, e.g., only "dry run")
\subsection{Case Study: OpenCode}
% same as above (not too deep/detailed, e.g., only "dry run")
% NOTE: Local model, hardware, runtime and settings are collected in the NOTE under Experimental Setup. Describe them once (there or here), and only refer to them in the other place.

% Note: detailed description of the case studies, except the logs/transcripts (usb-stick, cd, appendix). Address research questions with important moments in the case studies, e.g., when the agent got stuck, when it needed intervention, etc. 

% Note: Do not investigate full matrix: N apps x M approaches x K agents, but sample from it, with a focus on the case study with Claude Code (qualitative multi-case-study). Rest could be future work. 

% Note: Use consistent structure for the case studies